In Watt umgerechnet (FoTaBe S.~186) ist die mechanische Leistung $P = \Pm$. In der Zeit $t$ gibt der Motor die Arbeit $W = P\,t$ ab; dafür braucht er die Verbrennungswärme
\begin{align}
 Q &= \frac{W}{\eta} = \frac{P\,t}{\eta}
\end{align}
Die Treibstoffmasse folgt aus dem Heizwert von Heizöl, $H = \HuO$ (FoTaBe S.~193), $m = Q/H$, das Volumen aus der Dichte $\rho = \rhoHO$ (Heizöl EL, FoTaBe S.~187):
\begin{align}
 V &= \frac{m}{\rho} = \VolF \\
   &= \frac{\Pm\cdot\tS}{\etMO\cdot\Hu\cdot\rhoH} \\
   &= \Vol \approx \result{\VLP{0}{2}}
\end{align}
